\documentclass[11pt]{scrartcl}

\usepackage[stixtwo]{fontsetup}
\usepackage{amsmath}

\begin{document}
\begin{enumerate}
	\item \[
		      a_n
		      = \frac{\left(1 + \frac{3}{n}\right)^{n^2}}{\left(1 + \frac{3}{n^2}\right)^{n^2}}
		      \ge \frac{\left(\left(1 + \frac{3}{n}\right)^n\right)^n}{e^3}
		      \ge \frac{4^n}{e^3}
		      \to \infty,
	      \]
	      így a speciális rendőrelv miatt $a_n \to \infty$.

	\item \[
		      b_n
		      = \frac{
			      1 + \overbrace{\frac{1}{\sqrt{n}}}^{\to 0} + \overbrace{\frac{\sqrt[n]{n}}{n}}^{\to 0}
		      }{
			      3 + \underbrace{\sin(3n)}_{\text{korlátos}} \cdot \underbrace{\frac{1}{n}}_{\to 0}
		      }
		      \to \frac{1}{3}
	      \]

	\item Ha $n$ páros, akkor
	      \(
	      3 = \sqrt[n]{3^n} \le c_n = \sqrt[n]{2^n + 3^n} \le \sqrt[n]{3^n + 3^n}
	      = \sqrt[n]{2} \cdot 3 \to 3
	      \),
	      így ezekre $c_n \to 3$.

	      Ha $n$ páratlan, akkor
	      \(
	      2 = \sqrt[n]{2^n}
	      \le c_n = \sqrt[n]{2^n + \left(\frac{1}{3}\right)^n}
	      \le \sqrt[n]{2^n + 2^n} = \sqrt[n]{2} \cdot 2 \to 2
	      \),
	      így ezekre $c_n \to 2$.

	      $c_n$-nek két torlódási pontja van, így nincs határértéke.
\end{enumerate}
\end{document}
